\subsection{Suficiencia de \texttt{target\_windows\_per\_repetition}}

Usando las mismas 150 corridas, la Tabla~\ref{tab:target-windows} reporta,
para cada kernel, el mínimo de ventanas útiles logradas en cualquiera de
las diez repeticiones (\texttt{windows\_achieved\_min}) frente al objetivo
configurado (50 para CPU, 5 para GPU). En los 15 kernels del catálogo, el
mínimo observado supera el objetivo por al menos un orden de magnitud
(\texttt{npb\_mg}, el caso más ajustado del lado CPU, con 629 ventanas
mínimas frente a un objetivo de 50; \texttt{rodinia\_backprop}, el más
ajustado del lado GPU, con 140 frente a un objetivo de 5).

\begin{table}[H]
\centering
\caption{Ventanas útiles logradas (media y mínimo sobre 10 repeticiones) vs. objetivo configurado.}
\label{tab:target-windows}
\small
\begin{tabular}{@{}lrrr@{}}
\toprule
\textbf{Kernel} & \textbf{Objetivo} & \textbf{Media} & \textbf{Mínimo} \\
\midrule
\texttt{dgemm\_n2048} & 50 & 2464{,}7 & 2457 \\
\texttt{npb\_bt} & 50 & 24920{,}5 & 24828 \\
\texttt{npb\_cg} & 50 & 6692{,}7 & 6605 \\
\texttt{npb\_ft} & 50 & 3718{,}4 & 3649 \\
\texttt{npb\_lu} & 50 & 16347{,}1 & 16212 \\
\texttt{npb\_mg} & 50 & 637{,}1 & 629 \\
\texttt{npb\_sp} & 50 & 16758{,}7 & 16601 \\
\midrule
\texttt{gpu\_dgemm\_n4096} & 5 & 798{,}2 & 750 \\
\texttt{rodinia\_backprop} & 5 & 148{,}0 & 140 \\
\texttt{rodinia\_dwt2d} & 5 & 375{,}2 & 366 \\
\texttt{rodinia\_heartwall} & 5 & 1064{,}3 & 1058 \\
\texttt{rodinia\_hotspot} & 5 & 1231{,}6 & 1165 \\
\texttt{rodinia\_lavamd} & 5 & 980{,}7 & 922 \\
\texttt{rodinia\_lud} & 5 & 1564{,}2 & 1413 \\
\texttt{rodinia\_myocyte} & 5 & 1152{,}4 & 1113 \\
\bottomrule
\end{tabular}
\end{table}

\textbf{Conclusión (con el alcance correcto del chequeo)}: los datos
demuestran con claridad que \texttt{target\_windows\_per\_repetition}
\textbf{no condiciona el catálogo actual} -- en ningún caso observado
(150 corridas) estuvo cerca de ser la restricción activa, y el margen es
tan amplio (10--500$\times$ el objetivo contando lecturas crudas) que el
valor exacto configurado (50 o 5) no cambia qué corridas se aceptan hoy.
Eso es una afirmación distinta, y más débil, de decir que 50 o 5 son un
\emph{mínimo estadísticamente suficiente}: las ventanas contiguas de un
mismo kernel están correlacionadas entre sí (no son extracciones
independientes de una distribución), y para los kernels GPU, la
Sección~\ref{sec:gpu-interval-ns} muestra que varias "ventanas" NVML
consecutivas pueden compartir el mismo valor de sensor subyacente -- 5
lecturas GPU no equivalen necesariamente a 5 observaciones físicas
distintas. \texttt{target\_windows\_per\_repetition} debe entenderse,
entonces, como un \textbf{guardarraíl de duración mínima de la corrida}
(evita aceptar un kernel que termina casi de inmediato, sin tiempo
suficiente para que el detector de calentamiento y el resto del pipeline
operen con sentido) y no como una prueba de suficiencia estadística de la
muestra resultante. En la práctica, para el catálogo actual, ese
guardarraíl nunca es la restricción activa -- los seis casos ya
documentados con evidencia negativa de calentamiento (ARC-86) son
precisamente ejemplos de kernels cerca del límite en duración total,
aunque no en conteo de ventanas -- pero esto no debe leerse como una
validación de que el dataset resultante tiene suficiencia estadística
para ningún uso posterior (por ejemplo, entrenamiento de un modelo en
Fase 2): esa es una pregunta distinta, que requiere su propio análisis
de partición por corrida/carga, balance de clases y evaluación fuera de
muestra.
